\documentclass[11pt]{article}
\usepackage[T1]{fontenc}
\usepackage{textcomp}
\usepackage[margin=0.75in]{geometry}
\usepackage{amsmath,amssymb,amsthm}
\usepackage{array,booktabs}
\usepackage{hyperref}
\setlength{\parindent}{0pt}
\newtheorem{theorem}{Theorem}
\theoremstyle{definition}
\newtheorem{definition}{Definition}
\title{Flow Proof Artifact: Nat-mul-two-zero}
\author{Generated by \texttt{flow doc proof}}
\date{Flow Proof Book}
\begin{document}
\maketitle
Two times zero is zero for natural numbers.\\[0.5em]
\textbf{Source.} peano — https://en.wikipedia.org/wiki/Peano\_axioms\\[0.5em]
\bigskip
\noindent{\large \textbf{Derived fact 1.} \textit{2 * 0 = 0} \hfill the natural numbers \cdot multiplication \cdot two times zero is zero}\par
\smallskip\noindent\textit{Coordinate: the natural numbers \cdot multiplication \cdot two times zero is zero}\par
\smallskip\noindent\textit{Source: peano}\par
\smallskip\noindent\textit{Built on: zero is the right annihilator, for multiplication on the natural numbers}\par
\medskip
\noindent\textbf{Goal.} 2 * 0 = 0\par
\noindent$\displaystyle two \cdot 0 = 0$\par
\medskip
\noindent
\begin{tabular}{@{} >{\raggedright\arraybackslash}p{0.06\textwidth} >{\raggedright\arraybackslash}p{0.47\textwidth} >{\raggedleft\arraybackslash}p{0.41\textwidth} @{}}
\textbf{\#} & \textbf{Proof} & \textbf{Mathematics} \\
\midrule
\textbf{1.} & We prove that two times zero is zero for multiplication on the natural numbers. & \\[0.45em]
\textbf{2.} & Let two = succ(succ(0)). & \\[0.45em]
\textbf{3.} & We invoke the definitional clause governing multiplication on the natural numbers: zero is the right annihilator, for multiplication on the natural numbers (instantiated for two). & \\[0.45em]
\textbf{4.} & From step 2 and step 3, this implies two times 0 equals 0. Hence proven. & $\displaystyle two \cdot 0 = 0$ \\[0.45em]
\bottomrule
\end{tabular}
\label{thm:Nat-multiplication-two_times_zero_is_zero}
\par\medskip\hrule
\end{document}
